% Fourteen representative nodes; one physical template for all three timelines.
\begin{tikzpicture}[x=1mm,y=1mm,font=\figurefont\fontsize{8}{10}\selectfont,
 history card/.style={draw=FigureBlue!65,fill=white,rounded corners=2mm,line width=.65pt,minimum width=68mm,minimum height=18mm},
 history mark/.style={circle,minimum size=4mm,inner sep=0pt,draw=FigureBlue,fill=FigureBlueFill,line width=.7pt}]
\path[use as bounding box] (0,-5) rectangle (156,184);
\node[anchor=west,font=\figurefont\fontsize{8}{10}\selectfont,text=FigureInk] at (1,181) {经典模型：从参数估计到结构与组合};
\node[history card] at (34,169.0) {};
\node[anchor=west,font=\figurefont\fontsize{8}{10}\selectfont,text=FigureBlue] at (4,173.5) {1805—1809\quad 最小二乘};
\node[anchor=north west,align=left,text width=62mm,font=\FigureNoteFont\fontsize{7.5}{9.5}\selectfont,text=FigureInk] at (4,170.5) {平方残差组织参数拟合\\连接观测误差与优化};
\node[history card] at (122,169.0) {};
\node[anchor=west,font=\figurefont\fontsize{8}{10}\selectfont,text=FigureBlue] at (92,173.5) {1901\quad 主成分分析};
\node[anchor=north west,align=left,text width=62mm,font=\FigureNoteFont\fontsize{7.5}{9.5}\selectfont,text=FigureInk] at (92,170.5) {拟合低维线性子空间\\保留数据主要变化};
\node[history card] at (34,147.5) {};
\node[anchor=west,font=\figurefont\fontsize{8}{10}\selectfont,text=FigureBlue] at (4,152.0) {1936\quad 线性判别};
\node[anchor=north west,align=left,text width=62mm,font=\FigureNoteFont\fontsize{7.5}{9.5}\selectfont,text=FigureInk] at (4,149.0) {比较类间与类内散布\\学习分类投影方向};
\node[history card] at (122,147.5) {};
\node[anchor=west,font=\figurefont\fontsize{8}{10}\selectfont,text=FigureBlue] at (92,152.0) {1967\quad 近邻分类};
\node[anchor=north west,align=left,text width=62mm,font=\FigureNoteFont\fontsize{7.5}{9.5}\selectfont,text=FigureInk] at (92,149.0) {保留样本作为预测依据\\按距离形成局部判断};
\node[history card] at (34,126.0) {};
\node[anchor=west,font=\figurefont\fontsize{8}{10}\selectfont,text=FigureBlue] at (4,130.5) {1967\quad K-means};
\node[anchor=north west,align=left,text width=62mm,font=\FigureNoteFont\fontsize{7.5}{9.5}\selectfont,text=FigureInk] at (4,127.5) {用簇均值表示群组\\逐样本分配并更新均值};
\node[history card] at (122,126.0) {};
\node[anchor=west,font=\figurefont\fontsize{8}{10}\selectfont,text=FigureBlue] at (92,130.5) {1977\quad EM方法};
\node[anchor=north west,align=left,text width=62mm,font=\FigureNoteFont\fontsize{7.5}{9.5}\selectfont,text=FigureInk] at (92,127.5) {推断未观测变量\\按期望统计量更新参数};
\node[history card] at (34,104.5) {};
\node[anchor=west,font=\figurefont\fontsize{8}{10}\selectfont,text=FigureBlue] at (4,109.0) {1984\quad 分类与回归树};
\node[anchor=north west,align=left,text width=62mm,font=\FigureNoteFont\fontsize{7.5}{9.5}\selectfont,text=FigureInk] at (4,106.0) {递归划分输入区域\\在叶节点形成局部预测};
\node[history card] at (122,104.5) {};
\node[anchor=west,font=\figurefont\fontsize{8}{10}\selectfont,text=FigureBlue] at (92,109.0) {1995\quad 支持向量机};
\node[anchor=north west,align=left,text width=62mm,font=\FigureNoteFont\fontsize{7.5}{9.5}\selectfont,text=FigureInk] at (92,106.0) {以间隔选择分类边界\\用核内积支持非线性};
\node[history card] at (34,83.0) {};
\node[anchor=west,font=\figurefont\fontsize{8}{10}\selectfont,text=FigureBlue] at (4,87.5) {1995—1997\quad AdaBoost};
\node[anchor=north west,align=left,text width=62mm,font=\FigureNoteFont\fontsize{7.5}{9.5}\selectfont,text=FigureInk] at (4,84.5) {增加误分样本的权重\\加权组合弱分类器};
\node[history card] at (122,83.0) {};
\node[anchor=west,font=\figurefont\fontsize{8}{10}\selectfont,text=FigureBlue] at (92,87.5) {1996\quad DBSCAN};
\node[anchor=north west,align=left,text width=62mm,font=\FigureNoteFont\fontsize{7.5}{9.5}\selectfont,text=FigureInk] at (92,84.5) {连接密度可达样本\\发现非球形簇并标出噪声};
\node[history card] at (34,61.5) {};
\node[anchor=west,font=\figurefont\fontsize{8}{10}\selectfont,text=FigureBlue] at (4,66.0) {1996\quad Bagging};
\node[anchor=north west,align=left,text width=62mm,font=\FigureNoteFont\fontsize{7.5}{9.5}\selectfont,text=FigureInk] at (4,63.0) {重采样训练多个成员\\聚合不稳定预测};
\node[history card] at (122,61.5) {};
\node[anchor=west,font=\figurefont\fontsize{8}{10}\selectfont,text=FigureBlue] at (92,66.0) {2001\quad 随机森林};
\node[anchor=north west,align=left,text width=62mm,font=\FigureNoteFont\fontsize{7.5}{9.5}\selectfont,text=FigureInk] at (92,63.0) {扰动样本与特征候选\\降低树成员的相关性};
\node[history card] at (34,40.0) {};
\node[anchor=west,font=\figurefont\fontsize{8}{10}\selectfont,text=FigureBlue] at (4,44.5) {2001\quad 梯度提升};
\node[anchor=north west,align=left,text width=62mm,font=\FigureNoteFont\fontsize{7.5}{9.5}\selectfont,text=FigureInk] at (4,41.5) {逐轮拟合目标修正方向\\组合成加性预测模型};
\node[history card] at (122,40.0) {};
\node[anchor=west,font=\figurefont\fontsize{8}{10}\selectfont,text=FigureBlue] at (92,44.5) {2016—2017\quad XGBoost／LightGBM};
\node[anchor=north west,align=left,text width=62mm,font=\FigureNoteFont\fontsize{7.5}{9.5}\selectfont,text=FigureInk] at (92,41.5) {稀疏统计与近似分裂\\改进大规模树训练};
\begin{scope}[on background layer]
\draw[FigureMuted!35,line width=2.2mm,line cap=round] (73,169.0) -- (83,169.0) .. controls (91,164.0) and (65,152.5) .. (73,147.5) -- (83,147.5) .. controls (91,142.5) and (65,131.0) .. (73,126.0) -- (83,126.0) .. controls (91,121.0) and (65,109.5) .. (73,104.5) -- (83,104.5) .. controls (91,99.5) and (65,88.0) .. (73,83.0) -- (83,83.0) .. controls (91,78.0) and (65,66.5) .. (73,61.5) -- (83,61.5) .. controls (91,56.5) and (65,45.0) .. (73,40.0) -- (83,40.0);
\draw[white,dashed,line width=.3mm,dash pattern=on 1.5mm off 1.3mm] (73,169.0) -- (83,169.0) .. controls (91,164.0) and (65,152.5) .. (73,147.5) -- (83,147.5) .. controls (91,142.5) and (65,131.0) .. (73,126.0) -- (83,126.0) .. controls (91,121.0) and (65,109.5) .. (73,104.5) -- (83,104.5) .. controls (91,99.5) and (65,88.0) .. (73,83.0) -- (83,83.0) .. controls (91,78.0) and (65,66.5) .. (73,61.5) -- (83,61.5) .. controls (91,56.5) and (65,45.0) .. (73,40.0) -- (83,40.0);
\end{scope}
\node[history mark] at (73,169.0) {};
\draw[-{Stealth[length=1.4mm]},FigureStroke,line width=.6pt] (75.5,169.0) -- (80.5,169.0);
\node[history mark] at (83,169.0) {};
\node[history mark] at (73,147.5) {};
\draw[-{Stealth[length=1.4mm]},FigureStroke,line width=.6pt] (75.5,147.5) -- (80.5,147.5);
\node[history mark] at (83,147.5) {};
\node[history mark] at (73,126.0) {};
\draw[-{Stealth[length=1.4mm]},FigureStroke,line width=.6pt] (75.5,126.0) -- (80.5,126.0);
\node[history mark] at (83,126.0) {};
\node[history mark] at (73,104.5) {};
\draw[-{Stealth[length=1.4mm]},FigureStroke,line width=.6pt] (75.5,104.5) -- (80.5,104.5);
\node[history mark] at (83,104.5) {};
\node[history mark] at (73,83.0) {};
\draw[-{Stealth[length=1.4mm]},FigureStroke,line width=.6pt] (75.5,83.0) -- (80.5,83.0);
\node[history mark] at (83,83.0) {};
\node[history mark] at (73,61.5) {};
\draw[-{Stealth[length=1.4mm]},FigureStroke,line width=.6pt] (75.5,61.5) -- (80.5,61.5);
\node[history mark] at (83,61.5) {};
\node[history mark] at (73,40.0) {};
\draw[-{Stealth[length=1.4mm]},FigureStroke,line width=.6pt] (75.5,40.0) -- (80.5,40.0);
\node[history mark] at (83,40.0) {};
\node[anchor=west,font=\figurefont\fontsize{8}{10}\selectfont] at (1,27) {持续研究的问题};
\node[draw=FigureBlue,fill=FigureBlue!8!white,rounded corners=2mm,minimum width=49mm,minimum height=23mm,line width=.65pt] at (25,10) {};
\node[font=\figurefont\fontsize{8}{10}\selectfont] at (25,16) {表达与解释};
\node[align=center,font=\FigureNoteFont\fontsize{7.5}{9.5}\selectfont] at (25,8) {局部规则、特征交互\\可检查的预测依据};
\node[draw=FigureGreen,fill=FigureGreen!8!white,rounded corners=2mm,minimum width=49mm,minimum height=23mm,line width=.65pt] at (78,10) {};
\node[font=\figurefont\fontsize{8}{10}\selectfont] at (78,16) {训练与规模};
\node[align=center,font=\FigureNoteFont\fontsize{7.5}{9.5}\selectfont] at (78,8) {稀疏数据、近似搜索\\并行训练};
\node[draw=FigureYellow,fill=FigureYellow!8!white,rounded corners=2mm,minimum width=49mm,minimum height=23mm,line width=.65pt] at (131,10) {};
\node[font=\figurefont\fontsize{8}{10}\selectfont] at (131,16) {泛化与可靠性};
\node[align=center,font=\FigureNoteFont\fontsize{7.5}{9.5}\selectfont] at (131,8) {正则化、概率校准\\分布变化};
\end{tikzpicture}
